\documentclass[../../../include/open-logic-section]{subfiles}

\begin{document}

\olfileid{sth}{spine}{valpha}

\olsection{پړاوونه د $V_\alpha$ ګانو په توګه تعريفول}

په \olref[sth][ordinals][]{chap} کښې مو ښه ترتيبونه او د فون نويمن
ترتيبي عددونه تعريف کړل. په دې څپرکي کښې به دا د سټونو د مراتبو د لړۍ
\emph{پخپله} د مشخصولو دپاره وکاروو. د دې دپاره په ياد کړئ چې په
\olref[sth][ordinals][opps]{sec} کښې مو تالي او حدي ترتيبي عددونه تعريف
کړل. همدا مفهومونه په لاندې تعريف کښې کاروو:

\begin{defn}\ollabel{defValphas}
	\begin{align*}
	V_\emptyset &\defis \emptyset\\
	V_{\ordsucc{\alpha}} &\defis \Pow{V_\alpha} & &
	\text{د هر ترتيبي عدد دپاره }\alpha\\
	V_{\alpha} &\defis \bigcup_{\gamma < \alpha} V_\gamma & &
	\text{کله چې }\alpha\text{ حدي ترتيبي عدد وي}
\end{align*}
\end{defn}
\noindent
دا به په ترتيبي عددونو باندې د \emph{نامتناهيت هاخوا بازګښت} له لارې
تعريف وي. په دې اړه ئې بايد په طبيعي عددونو باندې د تابعو له بازګشتي
تعريفونو سره پرتله کړو.\footnote{په
\olref[sfr][infinite][dedekind]{sec} کښې د جمع، ضرب او د قوې تعريفونه
وګورئ.} د طبيعي عددونو په شان يو بنسټيز حالت او تالي حالتونه تعريفوو؛
خو د ترتيبي عددونو په کارولو کښې د \emph{حدي} حالتونو چلند هم بيانول
پکار دي.

د $V_\alpha$ ګانو دا تعريف به يو مهم پړاو وي. موږ د سټونو د مراتبو لړۍ
په غير رسمي توګه داسې توجيه کړې ده چې سټونه په \emph{پړاوونو} جوړېږي.
په حقيقت کښې $V_\alpha$ ګانې هماغه پړاوونه دي. خو دا د پړاوونو
\emph{دنننے} توصيف دے. د نظريې د کوم ممکن \emph{ماډل} د وړاندې کولو
پر ځای، پړاوونه به مو د سټونو د خپلې نظريې \emph{دننه} تعريف کړي وي.

\end{document}
